\documentclass[letterpaper,11pt]{article}

% Shared LaTeX preamble for CV and Cover Letter

\usepackage{iftex}

\ifPDFTeX
  \usepackage[T2A]{fontenc}
  \usepackage[utf8]{inputenc}
\else
  \usepackage{fontspec}
\fi

\usepackage[english, russian]{babel}

\usepackage{latexsym}
\usepackage[empty]{fullpage}
\usepackage{titlesec}
\usepackage{marvosym}
\usepackage[usenames,dvipsnames]{color}
\usepackage{verbatim}
\usepackage{enumitem}
\usepackage[hidelinks]{hyperref}
\usepackage{fancyhdr}
\usepackage{tabularx}
\usepackage{graphicx}

\input{glyphtounicode}
\pdfgentounicode=1

\pagestyle{fancy}
\fancyhf{} % clear all header and footer fields
\fancyfoot{}
\renewcommand{\headrulewidth}{0pt}
\renewcommand{\footrulewidth}{0pt}

\urlstyle{same}
\raggedbottom
\raggedright

\newcommand{\username}{alexo375}
\newcommand{\github}{https://github.com/\username}
\newcommand{\cpp}{C\texttt{++}}


% Adjust margins
\addtolength{\oddsidemargin}{-0.5in}
\addtolength{\evensidemargin}{-0.5in}
\addtolength{\textwidth}{1in}
\addtolength{\topmargin}{-.5in}
\addtolength{\textheight}{1.0in}

% Настройки абзацев для сопроводительного письма (без красной строки, с отступом между абзацами)
\setlength{\parindent}{0pt}
\setlength{\parskip}{1em}

\begin{document}

%----------ЗАГОЛОВОК----------
\begin{center}
  \begin{minipage}[b]{0.80\textwidth}
      \textbf{\Huge \scshape Горобченко Александр} \\ \vspace{4pt}
      \large Junior / Intern ML Engineer \\ \vspace{4pt}
      \small \href{tel:+79953755007}{+7 (995) 375-50-07} $|$ 
      \href{mailto:\username@yandex.ru}{\underline{\username@yandex.ru}} $|$ 
      \href{https://t.me/\username}{\underline{t.me/\username}} $|$
      \href{\github}{\underline{github.com/\username}} \\ \vspace{2pt}
      \small Рязань / Удаленная работа
  \end{minipage}%
  \begin{minipage}[b]{0.18\textwidth}
      \raggedleft
      \includegraphics[width=4cm, keepaspectratio]{me.png} 
  \end{minipage}
\end{center}

\vspace{0.5cm} % Отступ между шапкой и текстом письма
\hrule % Тонкая линия для визуального разделения
\vspace{0.5cm}

\textbf{Тема:} Отклик на стажировку: ML-инженер

Здравствуйте, коллеги!

Меня зовут Александр Горобченко, я студент магистратуры
Центрального Университета, изучаю ML с фокусом на компьютерное зрение.
Буду рад пройти конкурс на позицию стажера в Т-Банк.

Среди моих навыков --- уверенное владение Python (PyTorch, CatBoost, OpenCV) и C++,
а также опыт построения микросервисной архитектуры для ML-решений.
В роли R\&D инженера я разработал детерминированные и ML алгоритмы для решения задачи
обнаружения изменений на спутниковых изображениях площадной съемки высокого разрешения,
для чего с нуля собрал датасет, на котором обучил CNN-модель.
А в командном проекте, участвовавшем в акселерационной программе "Мой стартап 2025",
я создал микросервис для валидации изображений с использованием готовых ML моделей как
часть общего пайплайна приложения.
Помимо перечисленного мне также было бы интересно поработать с LLM или над задачами NLP.

Я хочу присоединиться к команде Т-Банка, так как сам являюсь клиентом,
постоянно использую ваши сервисы и вижу, насколько они технологичны.
Мне важно принимать участие в разработке продуктов, которые делают жизнь миллионов
людей удобной и комфортной.
Уверен, что мой опыт в CV и оптимизации алгоритмов поможет в автоматизации процессов
внутри экосистемы банка и принесет реальную пользу для улучшения клиентского опыта.

Стажировке готов уделять от 20 до 30 часов в неделю.

Буду рад познакомиться и более подробно рассказать о себе.

С уважением,\\
Александр Горобченко

\end{document}
